\documentclass[11pt,a4paper]{scrartcl}

\usepackage{amssymb}
\usepackage{amsmath}

\usepackage[utf8]{inputenc}
\usepackage[T1]{fontenc}
\newcommand{\ats}[1]{\par\noindent\textit{Atsakymas:} #1}
\usepackage{cancel}
\usepackage{tikz}
\usepackage{lmodern} 
\usepackage{amsmath,amssymb,amsthm}
\usepackage{graphicx}
\usepackage[dvipsnames,svgnames]{xcolor}
\usepackage{hyperref}
\usepackage{mathrsfs}
\usepackage[shortlabels]{enumitem}
\usepackage[obeyFinal,textsize=scriptsize,shadow,loadshadowlibrary]{todonotes}
\usepackage{mathtools}
\usepackage{microtype}

%%% Paragraph layout
\setlength{\parindent}{0pt}
\setlength{\parskip}{0.6\baselineskip}

%%% Colored sections
\renewcommand*{\sectionformat}%
  {\color{purple}\S\thesection\enskip}
\renewcommand*{\subsectionformat}%
  {\color{purple}\S\thesubsection\enskip}
\renewcommand*{\subsubsectionformat}%
  {\color{purple}\S\thesubsubsection\enskip}
\addtokomafont{paragraph}{\color{orange!35!black}\P\ }
\KOMAoptions{numbers=noenddot}

%%% Title
\addtokomafont{subtitle}{\Large}
\setkomafont{author}{\Large\scshape}
\setkomafont{date}{\Large\normalsize}
% Neigiamas vertikalus tarpas, kuris pakelia dokumento antraštę į viršų. 
% Galite keisti -2.5cm į -3cm (aukščiau) arba -1.5cm (žemiau).
\titlehead{\vspace*{-7.5cm}} 

% PRIDĖTA: Automatiškai perrašome datą į mokyklos pavadinimą
\AtBeginDocument{%
  \date{\normalsize Vilniaus licėjus}
}

%%% Page setup
\usepackage[headsepline]{scrlayer-scrpage}
\renewcommand{\headfont}{}
\addtolength{\textheight}{3.14cm}
\setlength{\footskip}{0.5in}
\setlength{\headsep}{10pt}
% Puslapio antraštė turi rodyti tik konkrečius dokumento duomenis.
% Nustatome juos aiškiai, kad neatsirastų "\@author" / "\@title" arba senų reikšmių.
\makeatletter
\AtBeginDocument{%
  \ihead{\footnotesize\textbf{Andrius Berniukevičius}}%
  \ohead{\footnotesize\textbf{\@title}}%
}
\makeatother
\automark{section}
\chead{}
\cfoot{\pagemark}

%%% Macros
\providecommand{\ol}{\overline}
\providecommand{\ul}{\underline}
\providecommand{\wt}{\widetilde}
\providecommand{\wh}{\widehat}
\providecommand{\eps}{\varepsilon}
\providecommand{\half}{\frac{1}{2}}
\providecommand{\inv}{^{-1}}
\newcommand{\dang}{\measuredangle} %% Directed angle
\newcommand{\vocab}[1]{\textbf{\color{blue}\sffamily #1}}
\providecommand{\CC}{\mathbb C}
\providecommand{\FF}{\mathbb F}
\providecommand{\NN}{\mathbb N}
\providecommand{\QQ}{\mathbb Q}
\providecommand{\RR}{\mathbb R}
\providecommand{\ZZ}{\mathbb Z}
\providecommand{\ts}{\textsuperscript}
\providecommand{\dg}{^\circ}
\providecommand{\ii}{\item}
\providecommand{\defeq}{\coloneqq}
\DeclareMathOperator*{\lcm}{lcm}
\DeclareMathOperator*{\argmin}{arg min}
\DeclareMathOperator*{\argmax}{arg max}
\providecommand{\hrulebar}{\par
\hspace{\fill}\rule{0.95\linewidth}{.7pt}\hspace{\fill}
\par\nointerlineskip \vspace{\baselineskip}}

%%% Optional colored theorems
\usepackage{thmtools}
\usepackage[framemethod=TikZ]{mdframed}
\mdfdefinestyle{mdbluebox}{%
  roundcorner=10pt,
  linewidth=1pt,
  skipabove=12pt,
  innerbottommargin=9pt,
  skipbelow=2pt,
  linecolor=blue,
  nobreak=true,
  backgroundcolor=TealBlue!5,
}
\declaretheoremstyle[
  headfont=\sffamily\bfseries\color{MidnightBlue},
  mdframed={style=mdbluebox},
  headpunct={\\[3pt]},
  postheadspace={0pt}
]{thmbluebox}
\mdfdefinestyle{mdredbox}{%
  linewidth=0.5pt,
  skipabove=12pt,
  frametitleaboveskip=5pt,
  frametitlebelowskip=0pt,
  skipbelow=2pt,
  frametitlefont=\bfseries,
  innertopmargin=4pt,
  innerbottommargin=8pt,
  nobreak=true,
  backgroundcolor=Salmon!5,
  linecolor=RawSienna,
}
\declaretheoremstyle[
  headfont=\bfseries\color{RawSienna},
  mdframed={style=mdredbox},
  headpunct={\\[3pt]},
  postheadspace={0pt},
]{thmredbox}
\mdfdefinestyle{mdgreenbox}{%
  skipabove=8pt,
  linewidth=2pt,
  rightline=false,
  leftline=true,
  topline=false,
  bottomline=false,
  linecolor=ForestGreen,
  backgroundcolor=ForestGreen!5,
}
\declaretheoremstyle[
  headfont=\bfseries\sffamily\color{ForestGreen!70!black},
  bodyfont=\normalfont,
  spaceabove=2pt,
  spacebelow=1pt,
  mdframed={style=mdgreenbox},
  headpunct={ --- },
]{thmgreenbox}
\mdfdefinestyle{mdblackbox}{%
  skipabove=8pt,
  linewidth=3pt,
  rightline=false,
  leftline=true,
  topline=false,
  bottomline=false,
  linecolor=black,
  backgroundcolor=RedViolet!5!gray!5,
}
\declaretheoremstyle[
  headfont=\bfseries,
  bodyfont=\normalfont\small,
  spaceabove=0pt,
  spacebelow=0pt,
  mdframed={style=mdblackbox}
]{thmblackbox}

%% Aplinkos su rėmeliais (thmtools)
\declaretheorem[style=thmbluebox,name=Teorema]{theorem}
\declaretheorem[style=thmbluebox,name=Lema]{lemma}
\declaretheorem[style=thmbluebox,name=Savybė]{property}
\declaretheorem[style=thmbluebox,name=Teiginys]{proposition}
\declaretheorem[style=thmbluebox,name=Išvada]{corollary}
\declaretheorem[style=thmbluebox,name=Teorema,numbered=no]{theorem*}
\declaretheorem[style=thmbluebox,name=Lema,numbered=no]{lemma*}
\declaretheorem[style=thmbluebox,name=Teiginys,numbered=no]{proposition*}
\declaretheorem[style=thmbluebox,name=Išvada,numbered=no]{corollary*}
\declaretheorem[style=thmgreenbox,name=Algoritmas]{algorithm}
\declaretheorem[style=thmgreenbox,name=Algoritmas,numbered=no]{algorithm*}
\declaretheorem[style=thmgreenbox,name=Teiginys]{claim}
\declaretheorem[style=thmgreenbox,name=Teiginys,numbered=no]{claim*}
\declaretheorem[style=thmredbox,name=Pavyzdys]{example}
\declaretheorem[style=thmredbox,name=Pavyzdys,numbered=no]{example*}
\declaretheorem[style=thmblackbox,name=Pastaba]{remark}
\declaretheorem[style=thmblackbox,name=Pastaba,numbered=no]{remark*}

%% Standartinės amsthm aplinkos (Apibrėžimai ir užduotys)
\theoremstyle{definition}
\ifcsname conjecture\endcsname\else\newtheorem{conjecture}{Hipotezė}\fi
\ifcsname definition\endcsname\else\newtheorem{definition}{Apibrėžimas}\fi
\ifcsname fact\endcsname\else\newtheorem{fact}{Faktas}\fi
\ifcsname answer\endcsname\else\newtheorem{answer}{Atsakymas}\fi
\ifcsname ques\endcsname\else\newtheorem{ques}{Klausimas}\fi
\ifcsname exercise\endcsname\else\newtheorem{exercise}{Užduotis}\fi
\ifcsname problem\endcsname\else\newtheorem{problem}{Uždavinys}[section]\fi
\ifcsname conjecture*\endcsname\else\newtheorem*{conjecture*}{Hipotezė}\fi
\ifcsname definition*\endcsname\else\newtheorem*{definition*}{Apibrėžimas}\fi
\ifcsname fact*\endcsname\else\newtheorem*{fact*}{Faktas}\fi
\ifcsname answer*\endcsname\else\newtheorem*{answer*}{Atsakymas}\fi
\ifcsname ques*\endcsname\else\newtheorem*{ques*}{Klausimas}\fi
\ifcsname exercise*\endcsname\else\newtheorem*{exercise*}{Užduotis}\fi
\ifcsname problem*\endcsname\else\newtheorem*{problem*}{Uždavinys}\fi

\newcounter{answerssection}
\newcommand{\answerssection}{%
  \setcounter{answerssection}{\value{section}}%
  \section{Atsakymai}%
  \setcounter{problem}{0}%
  \renewcommand{\theproblem}{\arabic{answerssection}.\arabic{problem}}%
}

%% GALUTINIS SPRENDIMAS PROOF APLINKAI
%% --- PRADŽIA: Įrodymo aplinkos sutvarkymas ---
\makeatletter

% 1. Užtikriname, kad žodis būtų "Įrodymas", net jei "babel" paketas bando jį perrašyti
\AtBeginDocument{%
  \renewcommand{\proofname}{Įrodymas}%
  \@ifpackageloaded{babel}{%
    \addto\captionslithuanian{\renewcommand{\proofname}{Įrodymas}}%
  }{}%
}

% 2. Priverstinai pašaliname scrartcl klasės bold šriftą ir paliekame TIK italic
% 3. Sujungiame su tuščiu kvadratėliu dešiniajame kampe.
\renewcommand{\qedsymbol}{\ensuremath{\Box}}
\renewenvironment{proof}[1][\proofname]{\par
  \pushQED{\qed}%
  \normalfont \topsep6\p@\@plus6\p@\relax
  \trivlist
  \item[\hskip\labelsep
        \normalfont\itshape % Priverčia naudoti tik pasvirą šriftą, kaip nuotraukoje
    #1\@addpunct{.}]\ignorespaces
}{%
  \popQED\endtrivlist\@endpefalse
}
\newenvironment{solution}{\par
  \normalfont \topsep6\p@\@plus6\p@\relax
  \trivlist
  \item[\hskip\labelsep
        \normalfont\itshape
    \textit{Sprendimas}\@addpunct{.}]\ignorespaces
}{%
  \endtrivlist\@endpefalse
}
\newcommand{\ans}[1]{\par\noindent\textit{Atsakymas:} #1}
\makeatother


\title{Skaičių teorija: Skaičių anatomija ir dalumo požymiai}
\author{Andrius Berniukevičius}

\newenvironment{solution}
    {\begin{list}{}{\setlength{\leftmargin}{10pt}\setlength{\rightmargin}{10pt}}\item[]}
    {\end{list}}

\begin{document}

\maketitle

\section{Skaičių anatomija: Kas slepiasi viduje?}

Praeitame užsiėmime išmokome naudoti vertikalų brūkšnį ($a \mid b$) ir algebriškai įrodinėti dalumą. Tačiau kartais uždaviniai kalba ne apie pačius skaičius, o apie jų \textbf{skaitmenis}. 

Mokykloje jūs tiesiog parašote skaičių, pavyzdžiui, $372$. Visi supranta, kad tai „trys šimtai septyniasdešimt du“. Tačiau olimpiadinėje matematikoje, jei norime surasti nežinomą dviženklį skaičių ir pažymime jo skaitmenis raidėmis $x$ ir $y$, mes \textbf{negalime} tiesiog parašyti $xy$. Kodėl? Nes algebroje $xy$ reiškia skaičių \textbf{sandaugą} ($x \cdot y$).

Kad atskirtume skaičiaus užrašymą skaitmenimis nuo daugybos, matematikai naudoja brūkšnį viršuje:
\[ \overline{xy} \]
Tai reiškia dviženklį skaičių, kurio dešimčių skaitmuo yra $x$, o vienetų – $y$.

Bet su brūkšniu matematinių veiksmų neatliksi. Norėdami skaičiuoti, mes turime skaičių „išardyti“ (išskleisti dešimtaine sistema):
\begin{itemize}
    \item Dviženklis: $\overline{ab} = 10a + b$
    \item Triženklis: $\overline{abc} = 100a + 10b + c$
    \item Keturaženklis: $\overline{abcd} = 1000a + 100b + 10c + d$
\end{itemize}
\textit{Svarbu atsiminti: pirmasis skaitmuo (pvz., $a$) niekada negali būti nulis, o patys skaitmenys yra sveikieji skaičiai nuo 0 iki 9.}

% ==========================================
% DALUMO POŽYMIAI KAIP TEOREMOS
% ==========================================
\section{Kodėl veikia dalumo požymiai?}

Mokykloje jūs išmokote taisykles: „skaičius dalijasi iš 3, jei jo skaitmenų suma dalijasi iš 3“. Bet ar kada nors susimąstėte, \textbf{kodėl} tai veikia? Dabar, kai mokate naudoti $a \mid b$ notaciją ir algebrines dalumo savybes, jūs patys galite šią taisyklę įrodyti!

\begin{center}
\fbox{
\begin{minipage}{0.9\textwidth}
\textbf{Dalumo iš 9 (ir iš 3) įrodymas}

Paimkime bet kokį triženklį skaičių $\overline{abc}$. Išskleiskime jį:
\[ \overline{abc} = 100a + 10b + c \]
Gudrybė: atskirkime tuos skaičius, kurie akivaizdžiai dalijasi iš 9. Žinome, kad $100 = 99 + 1$, o $10 = 9 + 1$. Pakeiskime mūsų išraišką:
\[ \overline{abc} = (99a + a) + (9b + b) + c \]
Sugrupuokime narius:
\[ \overline{abc} = (99a + 9b) + (a + b + c) \]
\[ \overline{abc} = 9(11a + b) + (a + b + c) \]

Pirmasis dėmuo $9(11a + b)$ akivaizdžiai dalijasi iš 9 (ir iš 3). Kad \textbf{visa suma} (mūsų pradinis skaičius) dalintųsi iš 9, pagal Sumos taisyklę privalo dalintis ir antrasis dėmuo, t.y. $(a + b + c)$. O kas yra $(a+b+c)$? Tai yra skaičiaus skaitmenų suma!
\end{minipage}
}
\end{center}

\subsection*{Olimpiadininkų ginklas: Dalumas iš 11}

Olimpiadose dalumas iš 9 pasitaiko retai, nes jis per daug žinomas. Tikrasis olimpiadininkų ginklas yra dalumo iš $11$ požymis. 

Pabandykime sukonstruoti jį taip pat, kaip darėme su devynetu. Paimkime keturaženklį $\overline{abcd}$. 
Mums reikia rasti $10, 100, 1000$ artimiausius kaimynus, kurie dalijasi iš 11.
Žinome, kad $11 \mid 11$, $11 \mid 99$, $11 \mid 1001$.
Išskleiskime:
\[ \overline{abcd} = 1000a + 100b + 10c + d \]
\[ = (1001a - a) + (99b + b) + (11c - c) + d \]
\[ = (1001a + 99b + 11c) - a + b - c + d \]

Pirmasis skliaustas idealiai dalijasi iš 11. Vadinasi, kad visas skaičius dalintųsi iš 11, turi dalintis ir likusi dalis:
\[ 11 \mid (-a + b - c + d) \quad \text{arba, pakeitus ženklus,} \quad 11 \mid (a - b + c - d) \]

\textbf{Taisyklė:} Skaičius dalijasi iš 11 tada ir tik tada, kai jo skaitmenų suma su kintančiais ženklais (plius, minus, plius, minus...) dalijasi iš 11. Dažnai užtenka, kad ši suma būtų lygi $0$ (nes $11 \mid 0$).

% ==========================================
% PAVYZDŽIAI
% ==========================================
\section{Pavyzdžiai: Nuo atradimo iki įrodymo}

\begin{example}[Klasikinis įrodymas]
Įrodykite, kad bet koks keturaženklis skaičius-palindromas (skaitomas vienodai iš abiejų pusių, pavyzdžiui, $4774$) visada dalijasi iš $11$.
\end{example}

\begin{solution}
\textbf{Sprendimas:}\\
Bet kokį tokį skaičių galime užrašyti kaip $\overline{abba}$ (kur $a \neq 0$).
Turime du būdus tai įrodyti.

\textit{1 būdas (Išskleidimas):} 
$\overline{abba} = 1000a + 100b + 10b + a = 1001a + 110b$.
Kadangi $11 \mid 1001$ ($1001 = 11 \cdot 91$) ir $11 \mid 110$ ($110 = 11 \cdot 10$), reiškinį galime užrašyti kaip $11(91a + 10b)$. Jis tobulai dalijasi iš $11$.

\textit{2 būdas (Požymio taikymas):} 
Pritaikome skaitmenų su kintančiais ženklais sumos taisyklę:
$a - b + b - a = 0$. Kadangi $11 \mid 0$, skaičius visada dalijasi iš 11. \hfill $\square$
\end{solution}


\begin{example}[Dviženklių magija]
Dviženklio skaičiaus skaitmenų suma lygi 12. Pakeitus šio skaičiaus skaitmenis vietomis, gaunamas skaičius, kuris yra $54$ didesnis už pradinį. Raskite pradinį skaičių.
\end{example}

\begin{solution}
\textbf{Sprendimas:}\\
Tegu pradinis skaičius yra $\overline{xy}$. Sąlygoje pasakyta, kad $x+y = 12$.
Apkeitę skaitmenis gauname skaičių $\overline{yx}$.
Užrašome sąlygą lygtimi:
\[ \overline{yx} - \overline{xy} = 54 \]
Dabar šį užrašą paverčiame algebra (išskleidžiame):
\[ (10y + x) - (10x + y) = 54 \]
\[ 9y - 9x = 54 \]
Padalinę iš 9, gauname labai paprastą lygtį:
\[ y - x = 6 \]
Dabar turime paprastą sistemą: suma $x+y = 12$, o skirtumas $y-x = 6$.
Sudėję šias dvi lygtis gauname $2y = 18 \implies y = 9$.
Tuomet $x = 12 - 9 = 3$.
Pradinis skaičius buvo $39$ (galime patikrinti: $93 - 39 = 54$). \hfill $\square$
\end{solution}


\begin{example}[Paslėpti skaitmenys]
Kokie skaitmenys $x$ ir $y$ gali būti įrašyti skaičiuje $\overline{5x7y}$, kad jis tobulai dalintųsi iš $45$?
\end{example}

\begin{solution}
\textbf{Sprendimas:}\\
Požymio dalumui iš 45 nėra. Bet $45 = 5 \cdot 9$. Kadangi 5 ir 9 neturi bendrų daliklių, skaičius dalinsis iš 45 tik tada, kai dalinsis ir iš 5, ir iš 9.

1) \textbf{Dalumas iš 5:} Skaičius dalijasi iš 5 tik tada, kai jo paskutinis skaitmuo yra $0$ arba $5$. Vadinasi, $y = 0$ arba $y = 5$.

2) \textbf{Dalumas iš 9:} Skaičiaus skaitmenų suma turi dalintis iš 9.
\begin{itemize}
    \item Jei $y = 0$: Skaičius yra $\overline{5x70}$. Skaitmenų suma: $5+x+7+0 = 12+x$. Koks skaitmuo $x$ padarys $(12+x)$ dalų iš 9? Tik $x=6$ (gausime 18). Skaičius: $5670$.
    \item Jei $y = 5$: Skaičius yra $\overline{5x75}$. Skaitmenų suma: $5+x+7+5 = 17+x$. Kad $(17+x)$ dalintųsi iš 9, $x$ privalo būti $1$ (gausime 18). Skaičius: $5175$.
\end{itemize}
Atsakymas: Skaičiai gali būti $5670$ ir $5175$. \hfill $\square$
\end{solution}

\vspace{1cm}

% ==========================================
% UŽDAVINIAI
% ==========================================
\newpage
\section{Uždaviniai savarankiškam sprendimui (18 iššūkių)}

Šiuose uždaviniuose, kur reikia ieškoti dviženklių ar triženklių skaičių, būtinai išskleiskite juos panaudodami dešimtainę anatomiją ($10a+b$).

\begin{enumerate}
    \renewcommand{\labelenumi}{\textbf{\theenumi.}}
    \item Įrodykite, kad bet kokio dviženklio skaičiaus ir to paties skaičiaus, parašyto atvirkščia skaitmenų tvarka, suma visada dalijasi iš $11$. (T.y., įrodykite, kad $11 \mid (\overline{ab} + \overline{ba})$).
    \item Įrodykite, kad bet kokio dviženklio skaičiaus ir to paties skaičiaus, parašyto atvirkščia tvarka, skirtumas visada dalijasi iš $9$.
    \item Įrodykite, kad skaičius $\overline{abcabc}$ (kuris gaunamas du kartus iš eilės parašius tą patį triženklį skaičių) visada dalijasi iš $7$, $11$ ir $13$.
    \item Ar tiesa, kad bet kaip sudėliojus devynis skaitmenis $1, 2, 3, 4, 5, 6, 7, 8, 9$, gautas devynženklis skaičius visada dalinsis iš $9$? Atsakymą pagrįskite.
    \item Įrašykite žvaigždutės vietoje tokį skaitmenį, kad skaičius $12345*$ dalytųsi iš: 
    a) $9$; \quad b) $8$; \quad c) $11$.
    \item Raskite visus galimus skaitmenis $x$ ir $y$, kad skaičius $\overline{7x3y}$ dalintųsi iš $15$.
    \item Raskite visus tokius dviženklius skaičius $\overline{ab}$, kurie yra $4$ kartus didesni už savo skaitmenų sumą.
    \item Prie dviženklio skaičiaus dešinėje pusėje prirašius $0$, jis padidėjo $252$ vienetais. Koks tai skaičius?
    \item Raskite trupmeną, kurios skaitiklis ir vardiklis yra dviženkliai skaičiai ir kuri nesikeičia, jei jos skaitiklio ir vardiklio skaitmenys sukeičiami vietomis. (T.y. raskite $\frac{\overline{ab}}{\overline{cd}} = \frac{\overline{ba}}{\overline{dc}}$, kai $\overline{ab} \neq \overline{cd}$).
    \item Ar įmanoma rasti tokį triženklį skaičių, kuris būtų lygiai $2$ kartus didesnis už savo paties išraišką atvirkščia skaitmenų tvarka? T.y., ar egzistuoja sprendinys lygčiai $\overline{abc} = 2 \cdot \overline{cba}$?
    \item Yra žinoma, kad penkiaženklis skaičius $\overline{24x8y}$ dalijasi iš $9$, o atėmus iš jo skaičių, užrašytą tais pačiais skaitmenimis atvirkščia tvarka, gaunamas skaičius, besidalijantis iš $11$. Raskite skaitmenis $x$ ir $y$.
    \item Įrodykite, kad lygtyje $\overline{ab} \cdot \overline{cd} = \overline{ba} \cdot \overline{dc}$ neįmanoma rasti tokių skaitmenų, kad visi jie būtų skirtingi nelyginiai skaičiai.
    \item Raskite mažiausią natūralųjį skaičių, kurio skaitmenų suma lygi $2026$.
    \item Raskite visus triženklius skaičius, kurie yra lygūs savo skaitmenų faktorialų sumai. T.y. $\overline{abc} = a! + b! + c!$ (Prisiminkite, kad $n! = 1 \cdot 2 \cdot \ldots \cdot n$, ir žinokite, kad $0! = 1$). \textit{Užuomina: koks didžiausias gali būti skaitmuo, kad jo faktorialas neviršytų 999?}
    \item Įrodykite dalumo iš $8$ požymį: „skaičius dalijasi iš $8$ tada ir tik tada, kai jo paskutinių trijų skaitmenų sudarytas skaičius dalijasi iš $8$“. \textit{Užuomina: atskirkite $1000$.}
    \item Turime skaičių, užrašytą iš lygiai $100$ vienetų: $111\dots111$. Ar jis dalijasi iš $11$? O iš $9$?
    \item Raskite visus keturaženklius skaičius $\overline{abcd}$, kurie yra lygūs tiksliam skaičiaus $\overline{ab}$ ir $\overline{cd}$ sumos kvadratui. T.y. $\overline{abcd} = (\overline{ab} + \overline{cd})^2$.
    \item \textbf{Iššūkis:} Parodykite, kad neegzistuoja joks dviženklis skaičius $\overline{xy}$, tenkinantis lygtį $\overline{xy} = x^2 + y^2$.
\end{enumerate}

\end{document}